\documentclass{article}
\usepackage{amsmath,amssymb,amsthm}
\usepackage{algorithm}
\usepackage{algpseudocode}
\usepackage{booktabs}
\usepackage{hyperref}
\usepackage{xcolor}
\newcommand{\E}{\mathbb{E}}
\newcommand{\R}{\mathbb{R}}
\newcommand{\norm}[1]{\left\|#1\right\|}
\newtheorem{assumption}{Assumption}
\newtheorem{theorem}{Theorem}
\newtheorem{lemma}{Lemma}
\newtheorem{corollary}{Corollary}
\begin{document}

\section*{Gradient Descent–Ascent for Convex–Concave Saddle‑Point Problems}

\section*{Mathematical Formulation}
Let $L:\mathcal{X}\times\mathcal{Y}\rightarrow\R$ be a differentiable function where
$\mathcal{X}\subseteq\R^{d_x}$ and $\mathcal{Y}\subseteq\R^{d_y}$ are convex compact sets.  
We consider the saddle‑point problem  

\[
\min_{x\in\mathcal{X}}\;\max_{y\in\mathcal{Y}} L(x,y).
\]

Denote the partial gradients by  

\[
\nabla_x L(x,y)\in\R^{d_x},\qquad \nabla_y L(x,y)\in\R^{d_y}.
\]

Given a step size $\eta>0$, the **Gradient Descent–Ascent (GDA)** iterates are  

\[
\boxed{
\begin{aligned}
x_{t+1} &= \Pi_{\mathcal{X}}\!\bigl(x_t-\eta\,\nabla_x L(x_t,y_t)\bigr),\\[2mm]
y_{t+1} &= \Pi_{\mathcal{Y}}\!\bigl(y_t+\eta\,\nabla_y L(x_t,y_t)\bigr),
\end{aligned}}
\tag{1}
\]
where $\Pi_{\mathcal{X}}$ and $\Pi_{\mathcal{Y}}$ denote Euclidean projections onto the feasible sets.

\section*{Pseudocode}
\begin{algorithm}[h]
\caption{Gradient Descent–Ascent (GDA)}
\begin{algorithmic}[1]
\Require Initial points $x_0\in\mathcal{X},\;y_0\in\mathcal{Y}$, step size $\eta>0$, maximum iteration $T$.
\For{$t=0,1,\dots,T-1$}
    \State Compute partial gradients 
    \[
    g^x_t \gets \nabla_x L(x_t,y_t),\qquad 
    g^y_t \gets \nabla_y L(x_t,y_t).
    \]
    \State Gradient descent on $x$ (ascent on $y$):
    \[
    \tilde{x}_{t+1}\gets x_t-\eta\,g^x_t,\qquad
    \tilde{y}_{t+1}\gets y_t+\eta\,g^y_t.
    \]
    \State Project back onto feasible sets:
    \[
    x_{t+1}\gets \Pi_{\mathcal{X}}(\tilde{x}_{t+1}),\qquad
    y_{t+1}\gets \Pi_{\mathcal{Y}}(\tilde{y}_{t+1}).
    \]
\EndFor
\State \textbf{Output:} $(x_T,y_T)$ (or averages $\bar{x}_T,\bar{y}_T$).
\end{algorithmic}
\end{algorithm}

\section*{Dry Run Example}
We illustrate (1) on a *pure* minimisation problem $f(w)=(w-3)^2$ (i.e. $L(x,y)=f(x)$, no $y$).  
Take $w_0=0$, step size $\eta=0.1$, and run three iterations.

\begin{table}[h]
\centering
\begin{tabular}{c c c c}
\toprule
Iteration $t$ & $w_t$ & $\nabla f(w_t)=2(w_t-3)$ & $f(w_t)$\\
\midrule
0 & $0$ & $2(0-3)=-6$ & $(0-3)^2=9$\\
1 & $0-0.1(-6)=0.6$ & $2(0.6-3)=-4.8$ & $(0.6-3)^2=5.76$\\
2 & $0.6-0.1(-4.8)=1.08$ & $2(1.08-3)=-3.84$ & $(1.08-3)^2=3.6864$\\
3 & $1.08-0.1(-3.84)=1.464$ & $2(1.464-3)=-3.072$ & $(1.464-3)^2=2.3660$\\
\bottomrule
\end{tabular}
\caption{Three iterations of gradient descent on $f(w)=(w-3)^2$.}
\end{table}

The objective value decreases from $9$ to $\approx2.37$ after three steps, illustrating the descent direction of (1).

\newpage
\section*{Assumptions}
\begin{assumption}[Convex–concave structure]\label{ass:convconc}
For every $y\in\mathcal{Y}$, $L(\cdot ,y)$ is convex on $\mathcal{X}$;  
for every $x\in\mathcal{X}$, $L(x,\cdot )$ is concave on $\mathcal{Y}$.
\end{assumption}

\begin{assumption}[Smoothness]\label{ass:smooth}
There exists $L>0$ such that for all $(x,y),(x',y')\in\mathcal{X}\times\mathcal{Y}$,
\[
\norm{\nabla_x L(x,y)-\nabla_x L(x',y')}\le L\norm{(x,y)-(x',y')},
\]
\[
\norm{\nabla_y L(x,y)-\nabla_y L(x',y')}\le L\norm{(x,y)-(x',y')}.
\]
\end{assumption}

\begin{assumption}[Bounded domain]\label{ass:bounded}
The feasible sets are bounded:
\[
D_x \triangleq \max_{x\in\mathcal{X}}\norm{x-x^*}<\infty,\qquad
D_y \triangleq \max_{y\in\mathcal{Y}}\norm{y-y^*}<\infty,
\]
where $(x^*,y^*)$ is a saddle point.
\end{assumption}

\begin{assumption}[Step size]\label{ass:stepsize}
The step size satisfies 
\[
0<\eta\le \frac{1}{2L}.
\]
\end{assumption}

\section*{Main Convergence Theorem}
\begin{theorem}[Primal–dual gap $O(1/T)$]\label{thm:convergence}
Let Assumptions \ref{ass:convconc}–\ref{ass:stepsize} hold and let $\{(x_t,y_t)\}_{t=0}^{T}$ be generated by (1).  
Define the averaged iterates
\[
\bar{x}_T\triangleq\frac{1}{T}\sum_{t=0}^{T-1}x_t,\qquad
\bar{y}_T\triangleq\frac{1}{T}\sum_{t=0}^{T-1}y_t .
\]
Then the primal–dual gap satisfies
\[
\underbrace{\max_{y\in\mathcal{Y}}L(\bar{x}_T,y)-\min_{x\in\mathcal{X}}L(x,\bar{y}_T)}_{\displaystyle G_T}
\;\le\; \frac{D_x^2+D_y^2}{\eta\,T}.
\tag{2}
\]
Consequently $G_T = O(1/T)$.
\end{theorem}

\section*{Theoretical Properties}
\begin{itemize}
\item \textbf{Stability.} The projection operators guarantee that all iterates stay inside the compact sets, preventing divergence even when the gradient is large.
\item \textbf{Step‑size sensitivity.} Condition \eqref{ass:stepsize} is sufficient for monotone decrease of the Lyapunov function $V_t\triangleq \norm{x_t-x^*}^2+\norm{y_t-y^*}^2$; larger $\eta$ may violate (2).
\item \textbf{Comparison.} Compared with pure gradient descent on a convex function, GDA obtains the same $O(1/T)$ rate for the saddle‑point gap, matching the optimal rate for first‑order methods under smoothness.
\item \textbf{Special cases.} If $L(x,y)=f(x)+\langle Kx,y\rangle-g(y)$ with $f,g$ convex and $K$ linear, (1) reduces to the primal–dual algorithm of Chambolle–Pock; the bound (2) recovers their $O(1/T)$ guarantee.
\end{itemize}

\section*{Preliminaries (Definitions and Lemmas)}
\begin{lemma}[Smoothness inequality]\label{lem:smooth}
If $L$ satisfies Assumption \ref{ass:smooth}, then for any $(x,y)$ and $(x',y')$,
\[
L(x',y')\le L(x,y)+\langle\nabla_x L(x,y),x'-x\rangle
+\langle\nabla_y L(x,y),y'-y\rangle
+\frac{L}{2}\norm{(x',y')-(x,y)}^2.
\tag{3}
\]
\end{lemma}
\begin{proof}
Apply the standard $L$‑smoothness inequality to the joint variable $z=(x,y)$.
\end{proof}

\begin{lemma}[Monotone variational inequality]\label{lem:vi}
For a convex–concave $L$, the saddle point $(x^*,y^*)$ satisfies
\[
\langle \nabla_x L(x^*,y^*),x-x^*\rangle
\ge 0,\qquad
\langle \nabla_y L(x^*,y^*),y-y^*\rangle
\le 0,\qquad\forall (x,y)\in\mathcal{X}\times\mathcal{Y}.
\tag{4}
\]
\end{lemma}
\begin{proof}
Convexity in $x$ gives $L(x,y^*)\ge L(x^*,y^*)+\langle\nabla_x L(x^*,y^*),x-x^*\rangle$; rearranging yields the first inequality. Concavity in $y$ yields the second.
\end{proof}

\section*{Main Proof (Step‑by‑Step Derivation)}
We prove Theorem \ref{thm:convergence}. All non‑trivial steps are labelled.

\begin{proof}
Let $(x^*,y^*)$ be a saddle point. Define the Lyapunov function  

\[
V_t \triangleq \norm{x_t-x^*}^2+\norm{y_t-y^*}^2 .
\tag{5}
\]

\paragraph{[Step 1] Expansion of $V_{t+1}$.}
Using the update rule (1) and non‑expansiveness of Euclidean projection,
\[
\begin{aligned}
\norm{x_{t+1}-x^*}^2 
&\le \norm{x_t-\eta\nabla_x L(x_t,y_t)-x^*}^2\\
&= \norm{x_t-x^*}^2-2\eta\langle\nabla_x L(x_t,y_t),x_t-x^*\rangle
+\eta^2\norm{\nabla_x L(x_t,y_t)}^2 .
\end{aligned}
\tag{6}
\]
Similarly,
\[
\begin{aligned}
\norm{y_{t+1}-y^*}^2 
&\le \norm{y_t+\eta\nabla_y L(x_t,y_t)-y^*}^2\\
&= \norm{y_t-y^*}^2+2\eta\langle\nabla_y L(x_t,y_t),y_t-y^*\rangle
+\eta^2\norm{\nabla_y L(x_t,y_t)}^2 .
\end{aligned}
\tag{7}
\]

\paragraph{[Step 2] Summation.}
Adding (6) and (7) yields
\[
\begin{aligned}
V_{t+1}
&\le V_t
-2\eta\Bigl[\langle\nabla_x L(x_t,y_t),x_t-x^*\rangle
-\langle\nabla_y L(x_t,y_t),y_t-y^*\rangle\Bigr]\\
&\qquad+\eta^2\bigl(\norm{\nabla_x L(x_t,y_t)}^2+\norm{\nabla_y L(x_t,y_t)}^2\bigr).
\end{aligned}
\tag{8}
\]

\paragraph{[Step 3] Bounding the gradient norm.}
By $L$‑smoothness (Assumption \ref{ass:smooth}) and the bounded domain (Assumption \ref{ass:bounded}),
\[
\begin{aligned}
\norm{\nabla_x L(x_t,y_t)}^2+\norm{\nabla_y L(x_t,y_t)}^2
&\le 2L^2 \norm{(x_t,y_t)-(x^*,y^*)}^2\\
&= 2L^2 V_t .
\end{aligned}
\tag{9}
\]

\paragraph{[Step 4] Plug (9) into (8).}
\[
V_{t+1}\le V_t-2\eta\underbrace{\Bigl[\langle\nabla_x L(x_t,y_t),x_t-x^*\rangle
-\langle\nabla_y L(x_t,y_t),y_t-y^*\rangle\Bigr]}_{\displaystyle \Delta_t}
+2\eta^2L^2 V_t .
\tag{10}
\]

\paragraph{[Step 5] Choose $\eta\le \frac{1}{2L}$.}
Since $\eta\le \frac{1}{2L}$, we have $2\eta^2L^2\le\eta L\le\frac{1}{2}$. Hence
\[
V_{t+1}\le V_t-\eta\Delta_t .
\tag{11}
\]

\paragraph{[Step 6] Relating $\Delta_t$ to the primal–dual gap.}
By convexity in $x$ and concavity in $y$ (Lemma \ref{lem:vi}),
\[
\begin{aligned}
\Delta_t
&= \langle\nabla_x L(x_t,y_t),x_t-x^*\rangle
-\langle\nabla_y L(x_t,y_t),y_t-y^*\rangle\\
&\ge L(x_t,y_t)-L(x^*,y_t)+L(x_t,y^*)-L(x_t,y_t)\\
&= L(x_t,y^*)-L(x^*,y_t).
\end{aligned}
\tag{12}
\]

\paragraph{[Step 7] Summation over $t=0,\dots,T-1$.}
Summing (11) and using (12) gives
\[
\begin{aligned}
\sum_{t=0}^{T-1}\eta\bigl[L(x_t,y^*)-L(x^*,y_t)\bigr]
&\le \sum_{t=0}^{T-1}\bigl(V_t-V_{t+1}\bigr)\\
&= V_0-V_T\le V_0 .
\end{aligned}
\tag{13}
\]

\paragraph{[Step 8] Divide by $\eta T$ and use convexity/concavity of $L$ for averages.}
Because $L(\cdot ,y^*)$ is convex,
\[
\frac{1}{T}\sum_{t=0}^{T-1}L(x_t,y^*)\ge L\!\left(\bar{x}_T,y^*\right),
\]
and because $L(x^*,\cdot )$ is concave,
\[
\frac{1}{T}\sum_{t=0}^{T-1}L(x^*,y_t)\le L\!\left(x^*,\bar{y}_T\right).
\]
Thus from (13):
\[
\eta\Bigl[L(\bar{x}_T,y^*)-L(x^*,\bar{y}_T)\Bigr]\le \frac{V_0}{T}.
\tag{14}
\]

\paragraph{[Step 9] Maximisation / minimisation over $y$ and $x$.}
By definition of the saddle‑point gap,
\[
\max_{y\in\mathcal{Y}}L(\bar{x}_T,y)-\min_{x\in\mathcal{X}}L(x,\bar{y}_T)
\le L(\bar{x}_T,y^*)-L(x^*,\bar{y}_T).
\tag{15}
\]
Combining (14) and (15) yields
\[
G_T\le \frac{V_0}{\eta T}.
\tag{16}
\]

\paragraph{[Step 10] Bounding $V_0$.}
Using the diameter definitions in Assumption \ref{ass:bounded},
\[
V_0=\norm{x_0-x^*}^2+\norm{y_0-y^*}^2\le D_x^2+D_y^2 .
\tag{17}
\]

\paragraph{[Step 11] Final bound.}
Insert (17) into (16):
\[
G_T\le \frac{D_x^2+D_y^2}{\eta T},
\]
which is exactly the claim \eqref{eq:2}. $\square$
\end{proof}

\section*{Rate Derivation (With Constants)}
The bound \eqref{eq:2} can be rewritten as
\[
G_T\le \frac{C}{T},\qquad C\triangleq\frac{D_x^2+D_y^2}{\eta}.
\]
If we set the maximal admissible step size $\eta=\frac{1}{2L}$, the constant becomes $C=2L\,(D_x^2+D_y^2)$. Hence the algorithm attains the optimal $O(1/T)$ rate for first‑order methods on smooth convex–concave saddle points.

\section*{Special Cases}
\begin{itemize}
\item \textbf{Pure minimisation.} When $L(x,y)=f(x)$ does not depend on $y$, the update for $y$ is immaterial and (1) reduces to standard gradient descent. The bound becomes $\frac{D_x^2}{\eta T}$, identical to the classical result for convex smooth functions.
\item \textbf{Bilinear games.} For $L(x,y)=x^\top A y$ with $A\in\R^{d_x\times d_y}$, $L$ is convex in $x$ and concave in $y$, and $L$‑smoothness constant equals $\|A\|_2$. The convergence rate (2) matches the known $O(1/T)$ rate for the extragradient method, confirming that GDA is optimal for this subclass as well.
\end{itemize}

\section*{Discussion}
The proof hinges on a simple Lyapunov decrease argument and the monotonicity property (4) of convex–concave functions. No variance or stochasticity appears; thus the same analysis carries over verbatim to the stochastic setting after taking expectations, provided the unbiasedness and bounded‑variance assumptions are added (as in SGD). The step‑size restriction $\eta\le 1/(2L)$ is tight for the plain GDA scheme; larger steps may cause oscillations, a phenomenon observed empirically in adversarial training.

\end{document}